%% ma: L11-B16-0057
%% lop: 11
%% chuong: 5
%% bai: 16
%% dang: 11.16.4
%% loai: DS
%% muc: [TH, TH, VD, VD]
%% dap_an: "ĐSĐĐ"
%% nguon: "tai-lieu-hoc-tap-mon-toan-lop-11-theo-chuong-trinh-giao-duc-pho-thong-2018.pdf (Lớp Toán Thầy Duy, Phạm Lê Duy, THPT Chu Văn An), Chương 3 Bài 2, Phần 2, Câu 1 (tr.375)"
%% kien_thuc: "Lớp 11 Bài 16 (giới hạn của hàm số)"
%% trang_thai: chinh-thuc
%% ghi_chu: "Đã duyệt 03/11/2026."
\begin{ex}
Cho $L=\lim\limits_{x\to-\infty}\left(\sqrt{x^2+ax+5}+x\right)$. Khi đó
\choiceTF
{\True $L=5$ khi $a=-10$.}
{$L>0$ khi $a>0$.}
{\True $L<0$ khi $a>0$.}
{\True $L=-1$ thì $a$ là một nghiệm của phương trình $x^2-3x+2=0$.}
\loigiai{\sloppy $\sqrt{x^2+ax+5}+x=\dfrac{ax+5}{\sqrt{x^2+ax+5}-x}$. Với $x<0$ chia tử và mẫu cho $x$: $\dfrac{a+\frac5x}{-\sqrt{1+\frac ax+\frac5{x^2}}-1}\to-\dfrac a2$. Vậy $L=-\dfrac a2$.
a) $a=-10\Rightarrow L=5$. Đúng.
b) $a>0\Rightarrow L<0$, không phải $L>0$. Sai.
c) $a>0\Rightarrow L<0$. Đúng.
d) $L=-1\Rightarrow a=2$ và $2^2-3\cdot2+2=0$. Đúng.}
\end{ex}
